\cotefile{57}
\scanpage{50}
% archive sheet number: 49
\uline{Cor. 4} Soit $S'\overset{f}{\unsure{\longrightarrow|\longleftarrow}}S$ un morphisme de
préschémas \uline{relt. connexe}. On suppose que
\margin{\rule[-6\baselineskip]{0.35pt}{6\baselineskip}}$f$ est \unsure{fid.|fidm} plat et de type fini, ou que
$S'$ et $S$ sont des spectres de corps. \deleted{Alors}
\deleted{\unsure{$k$}} $k',k$ (donc $k'/k$ est une extension
primaire). \deleted{Alors le \unsure{foncteur}} Soient
\deleted{$\unsure{H\longrightarrow H\times_S S'|H=H\times_S S'}$} $H,G$ des schémas en
groupes \unsure{comm.} \inserted{de type fini} sur $S$, $H$ \deleted{étant} rigide,
\inserted{soit $H'=H\times_S S'$, $G'=G\times_S S'$.}
Alors l'application canonique
\[
\operatorname{Hom}_{S\text{-}\mathrm{gr}}(H,G)\longrightarrow\operatorname{Hom}_{S'\text{-}\mathrm{gr}}(H',G')
\]
est bijective. En particulier, le
foncteur $H\unsure{\text{---}|\longrightarrow|\mapsto}H'$ de la catégorie
\margin{\rule[-3\baselineskip]{0.35pt}{3\baselineskip}}des schémas en gr. commutatifs \inserted{de type fini}
sur $S$ \deleted{\unsure{de}} rigides, dans la catégorie
\deleted{des schémas en groupes commutatifs}
analogue sur $S'$, est pleinement fidèle.
\scanpage{51}
\mspage{8}
% archive sheet number: 50
Il semble plus prudent, dans la définition de
% slanted
\margin{notion : \uline{très rigide}}la rigidité \inserted{sur un \uline{corps}}, d'exiger que \unsure{ses|les|tous les} sous-schém. fermés \inserted{\unsure{relt.|rel.!|rel.}} connexes
de $G$ ait la $=$ propriété de densité. Cela
signifie alors que $G$ a une suite de
composition (sur $k$ \uline{alg. clos}) \deleted{\unsure{tq}} dont les facteurs,
sont $\unsure{Z_{\ell}|Z\ell|Z/\ell}$, $G_m$, $\mathrm{Ab}$, --- et \deleted{cette} \deleted{condition} \inserted{notion}
\inserted{de \uline{très rigide}}, comme on voit facilement, se comporte
bien par limites inductives, est stable
par passage : \unsure{aux} \unsure{ss-schémas|ks-schémas} \unsure{fermés|fermé} simples
sur $k$.

\deleted{\uline{Proposition} $G$ rigide sur $S$ loc. noeth., $u:G\longrightarrow H$ un morphisme de $S$-groupes, où $H$ est un schém. en groupes \inserted{\unsure{comm.}} de type fini sur $S$. Soit $N=\operatorname{Ker}u$. Alors $N$ est universellement ouvert sur $S$.}

\deleted{On peut supposer $S=\operatorname{Spec} V$, $V$ anneau de valuation discrète.}

\uline{Remarque} Même si $A$ est un schém.
\margin{\rule[-3\baselineskip]{0.35pt}{3\baselineskip}\kern2pt\rule[-3\baselineskip]{0.35pt}{3\baselineskip}}abélien sur un anneau de valuation discrète $V$,
$B_1$ et $B_2$ deux sous-schémas \inserted{en groupes finis} de $A$ \deleted{\unsure{plats}} \inserted{\uline{étales}}
sur $V$, de degré $p$, ayant $=$ fibre
en \unsure{le} pt fermé $s$, il est possible qu'on ait $B_1\ne B_2$
\scanpage{52}
% archive sheet number: 51
Pour le voir, on part d'une courbe elliptique
\deleted{$B$} $C$ sur $S$, telle que dans \unsure{$C$|$C_s$} il n'y ait
pas de pt d'ordre $p$, mais $C(K)$
ayant un pt d'ordre $p$, donc $C(\deleted{\unsure{K|k}}V)$
a un pt d'ordre $p$, soit $\eta$. On
prend une courbe elliptique $C'$ sur $S$
contenant un sous-groupe isomorphe à $(\mathbb Z/p)\unsure{/S|_S}$,
d'où une section $\xi$ de $C'$. On
pose $A=C'\times_V C$, on y considère
les sections $\alpha=(\xi,e)$, $\beta=(\xi,\eta)$. Ces
sections engendrent des sous-groupes \inserted{$\Gamma_1,\Gamma_2$} de
$A$ isomorphes à $(\mathbb Z/p)_S$ [car
on a deux morphismes $(\mathbb Z/p)_S\unsure{\longrightarrow|\text{---}}A$,
qui sont des \unsure{is.m|imm.} car leur composé avec $\mathrm{pr}_1:A\longrightarrow C'$ l'est].
Ces groupes ont $=$ \deleted{\unsure{section}} fibre en $s$, mais non
en le point générique $s_1$. Ici,
le noyau de $A\longrightarrow A/\Gamma_1\times A/\Gamma_2$ \deleted{est}
n'est pas \deleted{\unsure{univers.}} ouvert sur $S=\operatorname{Spec} V$, c'est
$S\amalg_s(\mathbb Z/p)_{k(s)}$ [qui est non \unsure{ramifié} sur $S$].
Cela met des limites très strictes aux
améliorations éventuelles du th. de rigidité.
\scanpage{53}
\mspage{9}
% archive sheet number: 52
\uline{Proposition} Soient $k$ un corps, $k'$ une extension
\margin{\rule[-4\baselineskip]{0.35pt}{4\baselineskip}}\uline{primaire} de $k$, $G$ un groupe \inserted{\unsure{comm.}} de type fini sur
$k$, $G'=G\otimes_k k'$, $H'$ un ss-groupe fermé de
$G'$, tel que $\deleted{H'_{0}}\inserted{{H'}^{\,(0)}_e}$ (la composante connexe \inserted{de $e$}) \deleted{\unsure{de la section}})
soit rigide. Alors il existe une \deleted{ss-groupe}
extension \uline{finie} \uline{radicielle} $k_1$ de $k$, \deleted{telle que},
\deleted{\unsure{purement}} et un ss-groupe \inserted{\unsure{réduct't|réduitt} fermé} $H_1$ de $G_1=G\otimes_k k_1$,
tels que $H'_1=H'\otimes_k k_1$ \deleted{\unsure{soit isom. à}} \deleted{\unsure{soit}} \deleted{\unsure{égal à}} $H_1\otimes_{k_1}k'_1$
(où $k'_1=k'\otimes_k k_1$) \deleted{\unsure{et|ait}} \inserted{aient} $=$ ensemble sous-jacent.
% slanted
\margin{? N.B. Si on veut \uline{\unsure{majorer|imposer}}, $H'$ \unsure{par un} $H$ \unsure{défini} \unsure{on|par} le \unsure{pourra|passage} au \unsure{moins} pour une extension \unsure{finie} radicielle $k_1$}
\uline{Cor. 0} Si $k'/k$ est \deleted{\unsure{sép.}} primaire séparable, \deleted{alors $H'_1\supset H_1\otimes_{k_1}k'_1$.}

% in the left margin beside Cor. 0 and Corollaire 1
\diagram{\begin{tikzcd}[ampersand replacement=\&,column sep=small,row sep=small]
k' \arrow[r,dash]\arrow[d,dash] \& k'_1 \arrow[d,dash]\\
k\arrow[r,dash] \& k_1
\end{tikzcd}}

\uline{Corollaire 1} Il existe une \deleted{\unsure{sous-schéma}}
partie fermée $H$ de $G$ telle que
$|H'|={}$ image inverse de $H$

\uline{Corollaire 2} Si $k$ est parfait, il existe un
ss-groupe $H$ de $G$ tel que $H'\supset H\otimes_k k'$ et
$H'$ a $=$ ens. sous-jacent.

\uline{Corollaire 3} Soit \inserted{$T'\longrightarrow S$ de type fini surjectif},
\deleted{$T$ de type fini sur $S$,}
à fibres relt \deleted{\unsure{conn.}} irréductibles, $S'\longrightarrow S$
% slanted
\margin{\deleted{\unsure{hyp. plus gén.} \uline{\unsure{sans vérif.}}}}
\deleted{ét. est universellement un \unsure{hom.} \unsure{?} \unsure{quotient} topologique,}
$G$ \deleted{\unsure{$S$-groupe}} \inserted{schéma en groupes \unsure{comm.}} de type fini sur $S$,
$H'$ un ss-schém. en groupes de $G'=G\times_S S'$, \inserted{sur $S'$}, \unsure{on}
suppose : \unsure{a)} $H'$ ouvert sur $S'$ \unsure{(b)} les comp.
conn. des fibres de $H'$ aux pts gén. des fibres de $S'/S$ sont
\inserted{\unsure{rigides|rigide}}
\scanpage{54}
% archive sheet number: 53
Sous ces conditions, il existe une partie
\margin{\rule[-5\baselineskip]{0.35pt}{5\baselineskip}}\deleted{fermée} $H$ de $G$ dont l'image inverse dans
$G'$ soit $H'$. (Donc si le morphisme
$S'\longrightarrow S$ \deleted{\unsure{fid.} de $S$ \unsure{un} \unsure{morph.} top. quotient}
\inserted{est top. submersif}, \deleted{\unsure{$H$}}
p. ex. est propre, ou plat avec $S$ loc. noeth)
alors $H$ est une partie fermée de $G$.
En tout cas, c'est ``ensemblistement un
ss-groupe''.

% slanted
\margin{\uline{Cor. 4} Sous les conditions du Cor. 3, supposons que $S=\unsure{\operatorname{Spec}|\operatorname{Spm}}(k)$, \unsure{avec} $k$ \uline{parfait}, \deleted{\unsure{$S'$ rigide}} \deleted{\illegible{1}} (\unsure{bien entendu|ceci entraîne} les conditions en question). Alors il existe un ss-groupe fermé $H$ de $G$, \unsure{simple sur $k$|réduit sur $k$}, tel que \deleted{\unsure{$H'=H\otimes_k S'$|$H'=H\times_S S'$}} \inserted{$H'\supset H\otimes_k S'$} \deleted{et que} \inserted{et que} $H'$ ait $=$ ensemble ss-jacent que $H\otimes_k S'$.}

Le corollaire 3 résulte facilement du
corollaire 2, en se ramenant au cas où
$S=\operatorname{Spec}(k)$, et où on prend pour $k'$ le corps
\deleted{\unsure{du pt}} résiduel \unsure{des|du} pt gén. de $S'$.

Démonstration de la proposition. \deleted{On peut
supposer que} Appliquant le th. de
rigidité à $H'_0$, on voit qu'il
existe un ss-groupe $H_0$ \deleted{\unsure{du|de}} $G$ tel que
$H'_0=H_0\otimes_k k'$. Passant au quotient par
$H_0$, on est ramené au cas où $H'$ est
un groupe \uuline{fini}. Il est donc annulé
par un entier $n$, donc contenu
dans $\,{}_nG'=({}_nG)\otimes_k k'$. Cela nous ramène
au cas où $G$ est fini sur $k$. \deleted{\unsure{La}}
\scanpage{55}
\mspage{10}
% archive sheet number: 54
\deleted{\unsure{utilise donc le|utilise alors le|utilisé dans la}}

\deleted{\uline{Lemme} Soit $X$ fini sur un corps $k$, $k'$ une extension primaire de $k$, $\unsure{H'|X'}$ un sous-schém. (nécessairement \unsure{fermé|fini}) de $X'=X\otimes_k k'$. Alors l'application \unsure{est canonique}}

Comme $k'/k$ est primaire, \inserted{et $G/k$ fini}, l'application
$\deleted{\unsure{X'}}G'\longrightarrow G$ est bijective (mettre un
\uline{lemme}) donc en tant qu'ensemble
$H'$ est l'image inverse d'un ss-ensemble
\deleted{\unsure{de}} $H$ de $G$. \deleted{\unsure{qui est muni de la}}
\inserted{D'ailleurs}, \deleted{\unsure{puisque que}} par
\deleted{\unsure{de}} une extension radicielle finie convenable
$k_1$ de $k$, \deleted{\unsure{on trouve une} fois que}
\inserted{($G\otimes_k k'$, \unsure{d'id|d'id.} \unsure{sep.} séparable sur $k'$ (\unsure{sic|vif})) donc}
\deleted{que} $G_1\unsure{:|=}G\otimes_k k_1$ a des extensions résiduelles
séparables. \deleted{En} Si donc $H_1$ est l'image
inverse de $H$ dans $G_1$, muni de
la structure \uline{réduite} induite, $H_1$ est
un \uline{sous-schéma} \inserted{\uline{en groupes}} fermé de $G_1$.

De plus, $H_1\otimes_{\unsure{k'_1|k_1}}k'_1$ a $=$ espace
sous-jacent que $H'_1=H'\otimes_k k'_1$. On a
fini.
\scanpage{56}
% archive sheet number: 55
\uline{Remarque} L'hyp. sur $H'$ est certainement
vérifiée si $G_0$ est \deleted{strictement} très rigide
(p. ex. un tore \unsure{affin|forme|affine}, ou un schéma
abélien)

\deleted{\uline{Cor. 5} \inserted{Sous les conditions du cor. 4, i.e.} Supposons $H'$ plat sur $S'$, soit}

\uline{Proposition}

\margin{\rule[-13\baselineskip]{0.35pt}{13\baselineskip}}$S'$ de type fini \deleted{\unsure{réduit et}} \uline{relativement irréductible}
\deleted{\unsure{et}} sur $k$ \inserted{corps} \uline{parfait}, $\deleted{\unsure{H}}G$ schéma en groupes
comm. de type fini sur $k$, $G'=G\times_S S'$,
$H'$ un ss-schéma en groupes fermé
de $G'$ qui est \deleted{\unsure{plat}} \inserted{ouvert (p. ex. plat)} sur $S'$, et \deleted{dont}
la composante connexe ${H'}^{(0)}_{\eta}$ de la fibre
générique est \uline{rigide} (pour ce dernier
point, il suffit que $G_0$ soit \inserted{très} rigide).
Alors (i) il existe un sous-schéma en
groupes $H$ de $G$, \unsure{simple|lisse} sur $k$, tel que
$H\otimes_k \deleted{\unsure{G}}S'$ soit majoré par $H$ et ait $=$
ensemble sous-jacent. \deleted{Sous ces conditions,}
\deleted{\unsure{proposons}} \deleted{il} Considérons \inserted{alors} le morphisme
canonique $u:G\longrightarrow G/H=\deleted{\unsure{G_2 G}}G_1$, et le
morphisme correspondant \deleted{$G'=G\otimes_k S'\longrightarrow \unsure{G/H}$}
$u':G\otimes_k k'\unsure{\longrightarrow|\text{---}}G\otimes_k k'/H\otimes_k k'=\deleted{\unsure{G'_2}}G'_1$. (ii) \deleted{\unsure{Alors}} il
\scanpage{57}
\mspage{11}
% archive sheet number: 56
\margin{\rule[-4\baselineskip]{0.35pt}{4\baselineskip}}existe un unique ss-schéma en groupes \inserted{fermé}
$H'_1$ de $G'_1$, tel que $H'={u'}^{-1}(H'_1)$. \deleted{$H'_1$ est
\uline{fini}, \uline{plat}} Et $H'_1$ est \deleted{\unsure{aussi}} \uline{fini} \deleted{\unsure{et plat}} et \uline{radiciel}
sur $S'$ \deleted{[et le degré sur $S'$ est \unsure{donc}
une fonction de la \unsure{compos.} \unsure{loc.} $p$, cf.]}.
% slanted
\margin{\deleted{(iii) Pour que} Cela étant le degré \unsure{du groupe|de groupe} est une puissance de $p$, l'exposant \unsure{constructib.}}
(iii) Pour qu'il soit \uline{plat} sur $S'$, il f. et s. que $H'$ le soit \unsure{$\varnothing$|$\varphi$}.
\deleted{Si de plus \unsure{sur} $H'$ est \uline{prop.} sur $S'$, alors}

En d'autres termes, l'étude des
\deleted{\unsure{sous-schémas}} familles plates de ss-groupes
de $G$, paramétrées par $S'$, dont le
\unsure{mb|mbr} générique ait une comp. connexe
rigide (condition inutile si $G$ est
très rigide) est ramenée au cas \deleted{des}
\deleted{\unsure{sous-groupes}} familles plates, paramétrées
\unsure{sur|Ss|$S$,} de ss-groupes \uline{radiciels} de
degré fixé d'un groupe quotient
$G/H$. \deleted{\unsure{(où $H^{(0)}$ est lisse)}}

% slanted
\margin{\uline{Corollaire} Si \inserted{$H'$ est plat sur $S$ et si} une des fibres de $H'$ est simple, \uline{\unsure{ttles|ttes}} les fibres le sont, et $H'\simeq H\otimes_k S'$.}

\uline{Démonstration de la proposition}. \deleted{\unsure{(i)} \unsure{(4)}}
résulte de \deleted{\unsure{la pr.}} précédent cor. 4,
\deleted{\unsure{et}}\inserted{(ii)} de la proposition plus générale :
\scanpage{58}
% archive sheet number: 57
\deleted{\uline{Proposition}} \inserted{Lemme} $S$ \deleted{\unsure{schéma en groupes noeth}} \inserted{préschém. loc. noeth},
\margin{\rule[-2\baselineskip]{0.35pt}{2\baselineskip}}$G$, \inserted{$G_1$,} préschém. en groupes \inserted{de type fini} sur $S$, \deleted{\unsure{et}}
plats sur $S$, $u:G\unsure{\longrightarrow|\text{---}}G_1$ un morphisme \inserted{\unsure{de groupe|de groupes|à fibres}}
\uline{surjectif}, (donc nécessairement plat, \unsure{réf|cf.|surj.}),
% slanted
\margin{conditions de finitude inutile, si on suppose à priori $u$ plat \unsure{---}}$N$ son noyau (qui est \inserted{donc} plat sur $S$).
Alors il y a correspondance biunivoque
\margin{\rule[-10\baselineskip]{0.35pt}{10\baselineskip}}entre les sous-schémas fermés \inserted{$H$} de
$G$ qui sont invariants par $\deleted{\unsure{H}}N$, et les
sous schémas fermés \inserted{$H_1$} de \deleted{\unsure{$k$}} $G_1$, à $H_1$
correspondant $H=u^{-1}(H_1)$. Pour que
$H$ soit \deleted{propre} plat sur $S$, il f. et s. que
$H_1$ le soit. Pour que $H$ soit un
ss-groupe de $G$, il faut et il suffit que
$H_1$ soit un ss-groupe de $G_1$. Pour
que \inserted{de plus} $H_1\longrightarrow$ soit radiciel surjectif
sur $S$, il f. et s. que $H$ et $N$ aient
$=$ \deleted{\unsure{esp.}} \inserted{ens.} sous-jacent. \deleted{Alors} (Comme
$N\longrightarrow S$ \inserted{(donc $H\longrightarrow S$)} est \inserted{universelt} ouvert) c'est un morphisme
\uline{fini} radiciel \uline{surjectif}.

AQT par \uuline{théorie de descente}.
\scanpage{59}
% archive sheet number: 58
\hfill Albanese
\scanpage{60}
% archive sheet number: 59
\uline{Schéma d'Albanese}.

Soit $X$ \deleted{projectif et} \inserted{fid} plat \inserted{\unsure{q. cpt.|qu.-proj.}} sur $S$, \deleted{$\underline O_S\simeq f_*(\underline O_X)$ univt}

\deleted{On dit} Données équivalentes

\deleted{Pour un espace}

(i) $P$ fibré principal homogène [pour descente fid plate et \unsure{q. cpte|qu.-proj.}]
sur $S$, $S$-groupe $G$, et un $S$-morphisme $\unsure{\psi|u|\varphi}:X\longrightarrow P$

(\deleted{\unsure{ii}}i bis) Un $S$-groupe $G$, et une donnée de descente effective sur
$G_X=G\times_S X$ pour $X\longrightarrow S$, et la structure de \deleted{groupe} fibré
pr. hom. \inserted{à droite} sous $G_X$ sur $G_X$.

% left margin, next to (i bis)
\diagram{\begin{tikzcd}[ampersand replacement=\&,row sep=small,column sep=small]
P\arrow[d]\&P_X\arrow[l]\arrow[d]\\
S\&X\arrow[l]\arrow[ul,dotted,"\psi"']
\end{tikzcd}}

(ii) Un \deleted{\unsure{hom}} $S$-groupe $G$, et un $S$-hom $\varphi:X\times_S X$ satisfaisant
la condition \inserted{symbolique} $\varphi(x,y)\varphi(y,z)=\varphi(x,z)$ [impliquant $\varphi(x,x)=e$],
\deleted{\unsure{et}} et une condition d'effectivité, [autom. vérifiée
si $G$ est affine sur $S$, ou si $X$ est \unsure{qu. pr.} sur $S$
et $X\longrightarrow S$ admet loc. des quasi-sections plates finies,
p. ex. $X\longrightarrow S$ projectif].

% left margin, beside the cocycle condition
\diagram{\begin{tikzcd}[ampersand replacement=\&,column sep=small]
S \& X \& X\times_S X\arrow[l,shift left=1.3]\arrow[l,shift right=1.3] \& X\times_S X\times_S X\arrow[l,shift left=2]\arrow[l]\arrow[l,shift right=2]
\end{tikzcd}}

Supposons enfin que \deleted{\unsure{$G$}} $X$ soit \deleted{\unsure{projectif}} propre sur $S$
et tel que $\underline O_S\simeq f_*(\underline O_X)$ univt, alors on a équivalence
avec

(iii) Un $S$-groupe $G$, et un $S$-hom $\varphi:X\times_S X$ satisfaisant
seulement $\varphi(x,x)=e$, et une condition d'effectivité
[qui est autom. vérifiée si $G$ est affine/$S$, ou $G$ :
\unsure{qu. proj.}/$S$ et $X/\unsure{s|S|e}$ admet loc. des quasi-sections finies
plates, p. ex. $X$ projectif sur $S$].

% slanted
\margin{définir aussi en termes de $G$, $\unsure{\mathrm{Hom}|\mathrm{hom}}_{X/S}(G)={}$ ens des $(P,u)$ avec $G$ donné, d'où foncteur \deleted{\unsure{et}} et \unsure{faisceau|préfaisceau} de Weil, on peut dire [\unsure{par} ``\unsure{préfaisceau} de Weil''] $S'\longmapsto\underbrace{\mathrm{Hom}_{S'}(X_{S'},G_{S'})}_{\unsure{h^0|h}(X_{S'}/S',G_{S'})}/\unsure{\mathrm{Hom}^{\mathrm{cts}}|\mathrm{Hom}^{\mathrm{ct}}}$ et on passe au \uline{faisceau} associé \unsure{---|\ldots}}

Pour simplifier, nous supposons par la suite
$f:X\longrightarrow S$ \uline{projectif}, et $\underline O_S\simeq f_*(\underline O_X)$ univt. Pour
fibré principal sous variable $P$, \inserted{\unsure{ou|au} groupe variable $G$ \deleted{\unsure{propre}} \inserted{\unsure{comm.|proj.}} \unsure{proj.} sur $S$}, considérons
$\mathrm{Hom}_S(X,P)$ comme foncteur covariant en $P$, \inserted{soit $T_X(P)$.} On veut
le représenter. Tout d'abord, ce foncteur est \uline{pro-représentable}. Comme dans la catégorie des $P$, les \uline{lim}
\scanpage{61}
% archive sheet number: 60
finies existent, et que $T_X$ est exact à g., on est
ramené à vérifier que tt couple $(P,\unsure{\psi|u|\varphi}\in T_X(P))$
est \deleted{\unsure{majoré}} dominé par un ``couple minimal''.

Interprétons en effet le couple en question par
$\varphi:X\times_S X\longrightarrow G$ satisfaisant $\varphi(x,x)=e$, alors \deleted{\unsure{$G$}} les sous-groupes fermés $G_\alpha$ de $G$ tels que $\varphi$ se factorise par $G_\alpha\times_S G_\alpha$
\deleted{forment} un ens filtrant décroissant, et \unsure{because le noeth|because la noeth}
il y a un objet minimal, soit $G'$. Le couple
$(G',\varphi')$ correspondant est le minorant cherché.

\deleted{Montrons que le \unsure{groupe} pro-objet \inserted{\unsure{strict}} $(G_\alpha,\varphi_\alpha)$ qui représente $T_X$ est \unsure{ess. constant}.}
Nous allons
\deleted{montrer que} \inserted{mettre} le foncteur précédent \deleted{\unsure{est}} sous une
forme simple lorsque on y astreint $G$ à être un
schém. \uline{abélien} $A$, donc $A=\underline{\mathrm{Pic}}^0_{B/S}$, où $B=\widehat A=\underline{\mathrm{Pic}}^0_{\unsure{B|A}/S}$.
En effet, \deleted{alors} \deleted{soit $T_X(B)$ \unsure{l'alg de} \unsure{donne la} \unsure{fibre} \unsure{point}}
prendre pour simplifier $S'=X$, et faire le chgt de base \unsure{$S'\to S$|$S'$} 
\[
\begin{aligned}
\unsure{\pi_X|\Pi_X|T_X}(A)&\simeq\mathrm{Hom}'_S(X\times_S X,A)\simeq\mathrm{Hom}_{S'\text{-}\unsure{\mathrm{point\acute es}|\mathrm{point\acute e}}}(X_{S'},A_{S'})=\mathrm{Hom}_{S'\text{-}\mathrm{point}}(X_{S'},\underline{\mathrm{Pic}}^0_{B_{S'}/S'})\\
&\simeq\underbrace{\unsure{\mathscr B|\mathcal B}_{S'}(X_{S'},B_{S'})}_{\substack{\text{\uline{classes de corr.}}\\\text{\uline{divisorielles}}}}\simeq\mathrm{Hom}_{S'\text{-}\mathrm{point.}}(B_{S'},\underline{\mathrm{Pic}}_{X_{S'}/S'})=\underline{\mathrm{Hom}}_{S'\text{-}\unsure{\mathrm{gr}|\mathrm{point}}}(B_{S'},(\underline{\mathrm{Pic}}_{X/S})_{S'})\\
&\underset{\substack{\text{d'après le}\\\text{théorème de}\\\text{la rigidité}}}{\simeq}\mathrm{Hom}_{S\text{-groupes}}(B,\underline{\mathrm{Pic}}_{X/S})\qquad\text{donc}
\end{aligned}
\]
\[
\boxed{\pi_X(A)\simeq\mathrm{Hom}_{S\text{-gr}}(\widehat A,\underline{\mathrm{Pic}}_{X/S})}
\]
qui est un résultat du type de ``dualité de Cartier''.

Or supposons que $\underline{\mathrm{Pic}}_{X/S}$ contienne un sous-schéma
\deleted{\unsure{$B_{X/S}$}} \inserted{\unsure{disons|de s}} \inserted{l'un ss-groupe \unsure{sch.}} \uline{abélien} $\underline{\mathrm{Pic}}^0_{X/S}$, \unsure{Alors|dans} \unsure{utilisant|satisfaisant} \deleted{les} \deleted{\unsure{ss}}
``pts d'ordre fini''
\scanpage{62}
% archive sheet number: 61
on vérifie aisément que pour tt schém. abélien $B/S$, on
a une \uline{bijection}
\[
\mathrm{Hom}_{S\text{-gr}}(B,\underline{\mathrm{Pic}}_{X/S})\xleftarrow{\sim}\mathrm{Hom}_{S\text{-gr}}(B,\underline B_{X/S})
\]
d'où par \uline{dualité} \uline{des} V.A. une bijection
\[
\left\{\begin{aligned}
\mathrm{Hom}_{S\text{-gr}}(\widehat A,\underline{\mathrm{Pic}}_{X/S})&\xleftarrow{\sim}\mathrm{Hom}_{S\text{-gr}}(\underline A^0_{X/S},A)\\
\text{où}\qquad \underline A^0_{X/S}&=\widehat{\underline B}_{X/S}
\end{aligned}\right.
\]

Donc \uline{pour Un schém. abélien} $A$ \uline{sur} $S$, \deleted{\uline{\unsure{il existe|variable}}}
\uline{la donnée d'un fibré principal hom.} $P$ \uline{sous} $A$
\uline{et d'un $S$-morphisme} $\psi:X\longrightarrow P$ \uline{équivaut \unsure{can.|ess.} à}
\uline{la donnée d'un hom} \inserted{de \uline{$S$-gr}} $u:\underline A^0_{X/S}\longrightarrow A$. De façon
précise, il y a un fibré principal hom.
$\underline A^1_{X/S}$ sous $\underline A^0_{X/S}$, canonique, \deleted{tel que} \deleted{\unsure{tel que}}
\inserted{et un $S$-morphisme} \deleted{\unsure{unique}} $\psi_{X/S}:X\longrightarrow\underline A^1_{X/S}$, et la correspondance
\unsure{explicitée|explicite} plus haut s'obtient en prenant :
\[
\deleted{X} \quad P=\underline A^1_{X/S}\times_{\underline A^0_{X/S}}(A,u)\qquad \ldots\ldots\quad //
\]
On appelle
$\underline A^0_{X/S}$ et $\underline A^1_{X/S}$ resp. les schémas d'Albanese
\uline{homogènes} et \uline{inhomogènes} de $X/S$. Ils sont
en général non isom. \uline{On notera que}
$(\underline A^0_{X/S},\underline A^1_{X/S},\psi_{X/S})$ \uline{représente le foncteur} $T_X(P)$ par construction
\scanpage{63}
% archive sheet number: 62
On notera en $=$ temps que \uline{par construction}, \deleted{\unsure{la construction|le constructeur}}
la formation de $\underline A^0_{X/S},A^1_{X/S},\psi_{X/S}$ commute au
changement de base quelconque.

[N.B. Il y a peut-être lieu de définir
\unsure{axiomatiquemt|axiomatiquement} un \uline{schém. d'Albanese} de $X/S$ par
la condition de représenter le foncteur $T_X(P)$, $P$
pr. hom. \unsure{ss|en|les} schém. abélien, muni d'une $X$-trivialisation,
et par le fait que la $=$ chose soit valable pour
le changt de base. Pour l'\uline{existence}, il \uline{suffit}
mais il n'est pas nécessaire sans doute que
l'on puisse trouver un sous-schém. abélien
de $\underline{\mathrm{Pic}}_{X/S}$ ayant $\underline{\mathrm{Pic}}^0_{X/S}$ comme ens sous-jacent.
Ainsi, il est plausible \unsure{que 'un|que un|qu'un} schém. d'Albanese
pour $X/S$ existe si $S$ est le spectre d'un corps
$k$, \deleted{\unsure{et}} si $\underline{\mathrm{Pic}}^0_{X/k}$ n'est pas propre sur $k$
[plus grand sous-schém. abélien de $\underline{\mathrm{Pic}}^0_{X/k}$ ; il faut
prouver évidemment une propriété universelle de
ce dernier, pour tt changement de base \ldots\ldots{} ]].

\diagram{$\left\{\vcenter{\hbox{\begin{tikzcd}[ampersand replacement=\&,row sep=small,column sep=small]
X\arrow[r] \& \underline A^1_{X/S}\\
\underline{\mathrm{Pic}}^{\unsure{0|\bullet}}_{X/S}\&\underline{\mathrm{Pic}}^0_{\underline A^1/S}\simeq\underline{\mathrm{Pic}}^0_{\underline A^0/S}\arrow[l]\arrow[dl,"\sim"']\\
\underline B\arrow[u]\&
\end{tikzcd}}}\right.$}
\scanpage{64}
% archive sheet number: 63
\deleted{\uline{Question}} \inserted{N.B.} Si $\underline{\mathrm{Pic}}_{X/S}$ est simple sur $S$ aux pts de $\underline{\mathrm{Pic}}^0_{X/S}$, \deleted{\unsure{et si}} [ou ce
qui revient au $=$, le long de \uuline{la section} unité \ldots{}]
\deleted{la réunion des \unsure{composantes de} $S$, en \unsure{présence|propres}} \deleted{Alors} alors
$\underline{\mathrm{Pic}}^0_{X/S}$ est \uline{ouvert} dans $\underline{\mathrm{Pic}}_{X/S}$ [c'est \unsure{pour} tt \unsure{préschém.|préschéma} en
groupes loc. de type fini G sur $S$, car s'il est \uline{simple} sur $S$,
alors \deleted{\unsure{en}} comme chaque fois qu'on peut lui \unsure{imposer}
\unsure{connexion|convenablement} \unsure{par|que} les fibres par une section des \unsure{uns}
\unsure{composantes} \inserted{\uline{fibres géom} \deleted{\unsure{irréductibles}} \uline{loc. intègres}} \deleted{\unsure{irréductibles, les pts}},
$G^0$ est ouvert en $G$].
% margin note, linked by a bent arrow to the bracket
\margin{$\underline{\mathrm{IV}}~7.(6.6.\deleted{\unsure{bis|ter}})$}

Supposons le \uline{propre} sur $S$, donc c'est un schém. abélien
sur $S$. Alors (a) les fibres de $\underline{\mathrm{Pic}}_{X/S}$ sont simples
(b) \deleted{\unsure{$\Omega$}}$\unsure{Q|\mathrm O|\mathcal O}_X$ est coh. plat sur $S$ en \deleted{\unsure{codim}} \uline{dim $1$}
\inserted{i.e. $\underline{\omega}_{\underline{\mathrm{Pic}}_{X/S}}$ est loc. libre \ldots}

\deleted{[ \unsure{Réuni|Réun.|Réc.}.} Est-il vrai réciproquement que (a) et (b)
impliquent que $\underline{\mathrm{Pic}}_{X/S}$ est simple sur $S$ le long de
la section unité ??? Noter qu'en car. résiduelles $0$,
(a) est automatique, et \deleted{\unsure{je}} je pense que (b)
implique en effet (a), dans la forme plus générale :
Si $G$ est un schém en groupes sur $S$ \deleted{\unsure{en car. $0$}}, $S$ à
corps résiduels de car. $0$, et si \deleted{\unsure{$\Omega$}} $\underline\omega_{G/S}$ [faisceau
conormal à la section unité] est \uline{loc. libre}, \uuline{alors} $G$ est
simple sur $S$ le long de la section unité \ldots\ldots

De plus, il y a lieu de se demander si pour $X/S$
projectif et simple, et la car. $0$, les conditions
qu'on vient de dire sont satisfaites automatiquemt. On
en a une démonstration \uline{purement Algébrique}
via $\underline{\mathrm{Pic}}$ dans le cas où $S$ est \uline{réduit}. Dans le
cas $S$ arbitraire, je ne sais rien, \uline{\unsure{au secours !|au concours !|au contraire !}}
